\chapter{Introdução}

% Substitua as orientações abaixo pelo texto definitivo do trabalho.
Apresente o tema, delimite o objeto de estudo e forneça o contexto necessário
para que o leitor compreenda o problema investigado.

\section{Problema de pesquisa}

Formule a pergunta central que será respondida pelo trabalho. A pergunta deve
ser clara, específica e compatível com os dados e métodos disponíveis.

\section{Objetivos}

\subsection{Objetivo geral}

Declare, em uma frase, o principal resultado que a pesquisa pretende alcançar.

\subsection{Objetivos específicos}

Liste as etapas necessárias para alcançar o objetivo geral.

\begin{itemize}
    \item Primeiro objetivo específico;
    \item Segundo objetivo específico;
    \item Terceiro objetivo específico.
\end{itemize}

\section{Justificativa}

Explique a relevância acadêmica, social ou profissional do estudo e indique a
contribuição esperada da pesquisa.

\section{Estrutura do trabalho}

Descreva brevemente o conteúdo dos capítulos seguintes.

